\section{讨论与结论}
\label{sec:discussion}
\label{sec:conclusion}

分区架构中的原子移动连接了前后门层的放置决策。将原子留在纠缠区可省去往返存储区的阱转移，也会持续占用门位并经历空闲受激；将其回迁后，新的存储位又决定后续搬运的起点和轨迹可行性。跨层移动的代价取决于门位、驻留和存储位的共同选择。

本文提出GA-LK，通过遗传搜索联合选择双比特门位和原子的驻留或回迁方案，再为回迁原子匹配目标存储阱。每组完整方案均按当前层及后续门层的执行损失评价。衰减多层前瞻从各候选完成当前层后的原子位置、阱位占用和时钟出发，逐层计算非相邻交互之间的阱转移、受激和相干损失。

在统一物理模型下，GA-LK相较ZAC\cite{lin2025zac}与路由感知放置\cite{stade2025routingaware}提高了两个基准集的保真度几何均值，同时减少平均重排批次和时延。受控比较分别显示了遗传搜索与多层前瞻配置的收益。将门位、驻留和回迁位置置于共同的执行代价下求解，为减少分区中性原子电路的误差累积提供了编译优化方法。
